\section{Preliminaries}

\subsection{Notation}

Let $q\ge 3$ be prime and let $\F_q$ be the field of $q$ elements.  Vectors are column vectors unless stated otherwise.  For $x\in\F_q^n$, $\wt(x)$ denotes Hamming weight and $\supp(x)$ denotes support.  For $s\in\F_q^n$, define the extended secret
\[
    \tilde{s}=(-s_1,\ldots,-s_n,1)\in\F_q^{n+1}.
\]
For $c=(a,b)\in\F_q^n\times\F_q$, decryption-style inner products are written
\[
    c\cdot \tilde{s}=-a\transpose s+b.
\]
All probabilities are over the randomness of key generation, encryption, and noise unless otherwise specified.

\subsection{q-ary symmetric noise}

For $0\le \eta\le 1$, define the q-ary symmetric noise distribution $\chi_{q,\eta}$ over $\F_q$ by
\[
    e=0 \text{ with probability }1-\eta,
    \qquad
    e=u \text{ with probability }\eta/(q-1)\text{ for each }u\in\F_q^*.
\]
This is the natural q-ary analogue of Bernoulli noise for LPN.  It is also useful for coding-theoretic decoding because an error coordinate is either correct or replaced by a uniformly random nonzero field element.

\begin{lemma}[q-ary piling-up]
\label{lem:qary-piling}
Let $E_1,\ldots,E_B$ be independent q-ary symmetric errors with error probability $\eta$.  Let $\alpha_i\in\F_q^*$ be nonzero coefficients.  Then
\[
    E=\sum_{i=1}^B \alpha_iE_i
\]
is q-ary symmetric with error probability
\[
    \eta_B=\frac{q-1}{q}\left(1-\left(1-\frac{q\eta}{q-1}\right)^B\right).
\]
In particular, $\eta_B\le B\eta$ by the union bound.
\end{lemma}

\noindent The proof is included in Appendix~\ref{app:piling-up}.  The exact expression is useful for compaction parameters; the union bound is useful for simple correctness estimates.

\subsection{Sparse q-ary LPN}

Let $S_{n,k,q}$ be the distribution over $\F_q^n$ that samples a uniformly random support of size $k$ and assigns each supported position a uniformly random nonzero element of $\F_q^*$.  A sparse q-ary LPN sample with secret $s\in\F_q^n$ is
\[
    (a, a\transpose s+e),
    \qquad a\leftarrow S_{n,k,q},\quad e\leftarrow\chi_{q,\eta}.
\]

\begin{assumption}[Sparse q-ary LPN]
\label{ass:sparse-lpn}
For parameters $(n,k,q,\eta,M)$, the distribution of $M$ sparse q-ary LPN samples
\[
    \{(a_i,a_i\transpose s+e_i)\}_{i=1}^M
\]
with $s\leftarrow\F_q^n$, $a_i\leftarrow S_{n,k,q}$, and $e_i\leftarrow\chi_{q,\eta}$ is computationally indistinguishable from
\[
    \{(a_i,u_i)\}_{i=1}^M,
    \qquad u_i\leftarrow\F_q.
\]
\end{assumption}

Sparse LPN is a coding-theoretic noisy-decoding assumption.  Classical LPN and related noisy linear-equation assumptions are connected to learning parity with noise, random linear codes, and average-case coding hardness \cite{BlumKalaiWasserman2003,Alekhnovich2003,BerlekampMcEliecevanTilborg1978}.  Sparse-LPN variants have recently been used for homomorphic encryption \cite{CorriganGibbsHenzingerKalaiVaikuntanathan2024}.

\subsection{Dense q-ary LPN/code encryption}

The compaction layer uses a standard q-ary LPN/code primitive.  A dense q-ary LPN sample is
\[
    (A,Ar+e)
\]
for secret $r\in\F_q^{n_c}$, matrix $A\in\F_q^{m_c\times n_c}$, and error vector $e\in\F_q^{m_c}$ with independent $\chi_{q,\eta_c}$ coordinates.  Message encryption adds a codeword $\mathsf{Code}(x)$ to $Ar+e$.  In the simplest prototype, $\mathsf{Code}(x)=x\one$ is a repetition code.  A paper-grade version should replace repetition by a stronger q-ary error-correcting code.

\begin{assumption}[q-ary LPN/code pseudorandomness]
\label{ass:clpn}
For parameters $(n_c,m_c,q,\eta_c,M_c)$, polynomially many ciphertext blocks of the form
\[
    (A_j,A_jr+e_j+\mathsf{Code}(x_j))
\]
are semantically secure for messages $x_j\in\F_q$ under the q-ary LPN assumption and the chosen code/decoder.
\end{assumption}

This assumption is intentionally stated at the primitive level because the eventual library may instantiate $\mathsf{Code}$ by repetition, BCH-like q-ary codes, LDPC-like codes, or CRT-based code decompositions.

\subsection{Security target}

The target is secret-key IND-CPA security with public evaluation key.  The adversary may see the evaluation key, request polynomially many encryptions, choose two challenge messages, receive an encryption of one of them, and run arbitrary public homomorphic evaluation.  Since evaluation is deterministic and public, security follows from ciphertext and evaluation-key pseudorandomness under the stated assumptions.
